\newvideofile{Induktivitaet1}{Induktivität Teil 1}
\v{
	%Videoframe 
	\begin{frame}{Lernziele}
		\begin{Lernziele}{Induktion und Induktivität}
			Du kannst
			\begin{itemize}
				\item die Charakeristika der Induktivität beschreiben.
				\item die Induktivität über die Feldgrößen berechnen.
				\item die Induktivität über den magnetischen Widerstand berechnen.
			\end{itemize}
		\end{Lernziele}
		\speech{InduktivitaetLernziele}{1}{Nachdem wir im vorangegangenen Video mit der Induktion die Grundlagen für die Induktivität kennengelernt haben, wollen wir uns nun das auf diesem Prinzip basierende Bauteil der Indutivität näher ansehen, das in der Lage ist, magnetische Energie zu speichern. Im Anschluss betrachten wir die beiden Wege zur Berechnung einer Induktivität, die wir dann an einem einfachen Beispiel erproben.}
		\silence{2}
	\end{frame}
}

\section{Induktivität \label{Spule}}
\b{\begin{frame} \ftx{Induktivität}
	\f{trim=0 0.7cm 0 0.5cm,clip,width=0.35\textwidth}{width=0.4\textwidth}{Spule4}{{\bf Luftspule aus Lackdraht.} Luftspulen besitzen keinen magnetischen Kern und haben dadurch eine vergleichsweise geringere Induktivität als Spulen mit weichmagnetischem Kern. }
	\end{frame}
}
\begin{frame} \ftx{Induktivität} \index{Induktivität}
	\s{
	 Die Induktivität beschreibt die Fähigkeit, elektrische Energie in einem magnetischen Feld zu speichern. Bauelemente, die diese Fähigkeit besitzen, werden Spule (Abbildung \ref{InduktivitaetB}), Drossel oder auch Induktivität genannt.
	}
	\b{
		\begin{columns}
		\column[c]{0.5\textwidth}
	}

	\fo{width=0.5\textwidth, alt={Ein Draht, der eine Schleife bildet. Durch die Leiterschleife fließt elektrischer Strom (I), wodurch ein Magnetfeld, dargestellt durch kreisförmige Linien um die Schleife herum, erzeugt wird. Die Polarität des Magnetfeldes wird durch die Buchstaben 'N' (Nordpol) und 'S' (Südpol) angezeigt. Die Richtung der Magnetfeldlinien wird durch Pfeile entlang der Linien verdeutlicht.}}{width=\textwidth}{Magnetisches_Feld_Leiterschleife_Induktion}{
		\put(31.5, 18){\color{current}$I$}
		\put(65, 4){\color{current}$I$}
	}{{\bf Elektromagnetisches einer Induktivität.} Die Induktivität beschreibt, in welchem Maße sich ein Magnetfeldum eine stromdurchflossene Leiterschleife aufbaut. \label{InduktivitaetB}}
	\s{
		Die Funktionsweise einer Induktivität beruht auf dem elektromagnetischen Feld, das durch den Stromfluss erzeugt wird. Dieses elektromagnetische Feld induziert einen Strom in benachbarten Leiterspulen des Bauelements. Durch diesen Effekt kann elektrische Energie vorübergehend in einem magnetischen Feld gespeichert werden. Zum besseren Verständnis wird im Folgenden die Funktion der Induktivität schrittweise erklärt.

		Eine Leiterschleife wird von Strom durchflossen. Dieser Strom erzeugt ein Magnetfeld, dessen Stärke proportional zur Stromstärke ist -- je stärker der Strom, desto stärker das Magnetfeld. Wird die Stromstärke kontinuierlich geändert, erzeugt dies ebenso eine Änderung des Magnetfeldes, welche wiederum (wie in Kapitel \ref{KapInduktion} beschrieben) eine Spannung in benachbarten Leiterschleifen induziert. Diese induzierte Spannung wirkt gemäß der Lenzschen Regel der Änderung des Stroms entgegen. Folglich wirkt beim Einschalten der Stromquelle an der Spule die induzierte Spannung dem Stromfluss entgegen, sodass der Strom erst allmählich ansteigt. Beim Ausschalten der Spule bewirkt die induzierte Spannung, dass der Strom noch eine Zeit lang nachfließt. Diese Gegebenheit, auch Selbstinduktion genannt, ermöglicht eine magnetische Speicherung.
		Die Induktivität gibt folglich an, wie stark sich ein elektrisches Bauteil selbstinduziert. Die Höhe der Induktität $L$ hängt ab von der Anzahl der Windung $N$, der Permeabilität des Spulenkerns $\mu_r$, der Querschnittsfläche des Kerns und der geometrischen Form des Bauteils. Für Induktivitäten gilt im Allgemeinen, dass das magnetische Feld sich proportional zum Wert des Stroms verhält.

		Rechnerisch lässt sich die Induktivität $L$ mit der Einheit Henry (H) auf zwei unterschiedliche Weisen ermitteln. Zum einen kann sie mittels der Relation des entstehenden Magnetfelds $N \cdot \varPhi$ zu dem verursachenden Strom $I$ (wie in Gleichung \ref{InduktivitaetFormel}) berechnet werden.
	}
	\pause

	\b{\column[c]{0.5\textwidth}}

	\begin{eq}
		L = \frac{N \cdot \varPhi}{I} \label{InduktivitaetFormel} \qquad [\mathrm{H}]
	\end{eq}

	\b{\pause}
	\s{
		Zum anderen kann die Induktivität bei einfachen geometrischen Strukturen, wie zum Beispiel zylindrische, rechteckige oder Ringkerspulen, über die Relation der Windungen $N^2$ zu dem magnetischen Widerstand \( R_{\mathrm{m}} \) bestimmt werden (Gleichung  \ref{InduktivitaetFormel2}).

		\begin{eq}
			L=\frac{N^2}{R_{\mathrm{m}}} \label{InduktivitaetFormel2} \qquad [\mathrm{H}]
		\end{eq}
		\f{width=0.5\textwidth, alt={Fotografie einer Ringkernspule. Der ringförmige Eisenkern ist von Kuperdraht umwickelt. Am unteren Ende befinden sich die Anschlüsse des Kupferdrahts.}}{height=0.7\textheight}{Spule_Ringkern1b}{{\bf Foto einer Ringkernspule.} \index{Ringkernspule}Die ringförmige Bauweise bringt die Besonderheit mit sich, dass sich der magnetische Fluss vor allem innerhalb des Kerns bewegt, wodurch magnetische Störfelder reduziert werden.}
	}
	\s{
		\begin{Merksatz}{Induktivität}
			Die Induktivität $L$ beschreibt ein Bauelement, das in der Lage ist, elektrische Energie in einem magnetischen Feld zu speichern.
		\end{Merksatz}
	}

	\s{
		Die Berechnung der Spannung einer Induktivität lässt sich von der Berechnung der Spannung einer Induktion ableiten (Gleichung \ref{GlInduktionsgesetz}).

		\begin{eq}
			u_{\mathrm{i}} = -N\cdot \frac{\d \varPhi}{\d t} \tag{\ref{GlInduktionsgesetz}}
		\end{eq}

		Die Spannung einer Induktivität $u$ ist eine zeitliche Änderung des Stromes ${\d i}/{\d t}$ multipliziert mit der Induktivität $L$. Auch hier findet die Lenzsche Regel Anwendung, welche ein negatives Vorzeichen zur Folge hat.
	}

	\begin{eq}
		u =-L\cdot\frac{\d i}{\d t}\label{GLInduktivitaet2}
	\end{eq}

	\s{
		In Schaltplänen werden Induktivitäten mittels schwarz gefüllter Rechtecke, wie in der oberen der beiden Abbildungen \ref{AbbSchaltzeichenInduktivitaet} dargestellt. Seltener wird ein gewundener Leiter, wie in der unteren Abbildung \ref{AbbSchaltzeichenInduktivitaet} für die Darstellungen einer Induktivität genutzt.
	}
	\pause
	\fu{
		\input{Tikz/src/schaltzeichen_induktivitaet}\alttext{Das Bild zeigt zwei elektrische Schaltkreise: Der obere Abschnitt zeigt einen rechteckigen, schwarz gefüllten Block, durch den ein Strom (I) fließt, während eine Spannung (U) angelegt ist. Die Pfeile zeigen die Richtung des Stroms und der Spannung an. Es ist das Schaltzeichen einer Spule nach nach IEC 617-4. Der untere Abschnitt zeigt eine Wellenlinie, ebenfalls mit der gleichen Strom- (i) und Spannungsrichtung (U). Es ist das Schaltzeichen einer Spule nach DIN EN 60617-4.}
	}{Schaltzeichen der Induktivität \label{AbbSchaltzeichenInduktivitaet}}

	\b{\end{columns}}
	\speech{InduktivitaetEinleitung}{1}{Wir haben bereits die wichtigsten magnetischen Größen kennengelernt – den magnetischen Fluss, die Flussdichte und die Durchflutung. Zusätzlich haben wir ihre Wirkungen betrachtet, zum Beispiel im Induktionsgesetz und bei der Lorentzkraft. Jetzt wird es Zeit, uns ein paar magnetische Bauteile anzuschauen. In der Elektrotechnik kennen wir bereits den Kondensator, der elektrische Ladungen speichert. Diese Ladungen werden im elektrischen Feld gespeichert. Wir hatten schon mehrfach die Parallelen zwischen elektrischen und magnetischen Größen angesprochen – zum Beispiel Spannung und Strom auf der einen Seite, und ihre magnetischen Entsprechungen auf der anderen. Daher ist es nur logisch, dass es auch ein magnetisches Bauteil gibt, das Energie speichern kann – ähnlich wie der Kondensator, nur eben im magnetischen Feld. Dieses Bauteil nennt man Induktivität, häufig auch Spule oder Drossel genannt. Die Induktivität ist also das magnetische Gegenstück zum Kondensator. Jedoch speichert die Induktivität Energie im magnetischen Feld.}
	\silence{2}
	\speech{FunktionsweiseInduktivitaet}{1}{Schauen wir uns nun das Funktions-prinzip einer Induktivität an. Dazu betrachten wir eine einfache Leiterschleife, durch die ein Strom fließt. Der Strom verläuft also einmal im Kreis durch den Leiter. Die Richtung des Stroms spielt dabei eine entscheidende Rolle. Mit der Rechten-Hand-Regel lässt sich bestimmen, wie das Magnetfeld um den Leiter verläuft. Um den Leiter entsteht also ein Magnetfeld, das einen magnetischen Fluss durch die Leiterschleife erzeugt. Die magnetische Flussdichte ist dabei der Fluss bezogen auf die Fläche der Schleife. In Analogie zum Kondensator können wir das Verhältnis zwischen dem gesamten magnetischen Fluss und dem verursachenden Strom betrachten. Dieses Verhältnis nennt man Induktivität. Sie beschreibt also, wie stark ein Stromfluss in einer Leiterschleife ein Magnetfeld aufbaut.} 
	\speech{FunktionsweiseInduktivitaet}{2}{Dieser Zusammenhang wird in der Formel deutlich: Das Produkt aus Windungsanzahl und magnetischem Fluss, geteilt durch die Stromstärke ist gleich der Induktivität. Die Induktivität wird mit dem Buchstaben L bezeichnet und in der Einheit Henry angegeben – benannt nach dem Physiker Joseph Henry. In der Theorie lässt sich die Induktivität mit dieser Formel berechnen. In der Praxis ist das jedoch eher ungenau, weil reale Spulen keine idealen Bedingungen erfüllen. Daher wird die Induktivität in der Praxis meistens durch Messung bestimmt.}
	\silence{2}
	\speech{FunktionsweiseInduktivitaet}{3}{Die Wirkung einer Induktivität lässt sich auch über Strom und Spannung beschreiben. An einer Induktivität liegt eine Spannung an, die gleich minus L mal D-I durch D-T ist. Das erinnert stark an das Induktionsgesetz, und tatsächlich besteht hier ein enger Zusammenhang. Die Spannung an der Induktivität ist also proportional zur zeitlichen Änderung des Stroms, multipliziert mit dem Wert der Induktivität L. Auch hier gilt die Lenzsche Regel: Ursache und Wirkung sind entgegengesetzt. Das Minuszeichen in der Formel zeigt genau diesen Zusammenhang an.}
	\silence{2}
	\speech{FunktionsweiseInduktivitaet}{4}{Das Schaltzeichen einer Induktivität sieht dem Schaltzeichen des Widerstands ähnlich. In der europäischen Norm wird es als schwarzes, ausgefülltes Rechteck dargestellt. Manchmal wird euch auch eine andere Variante begegnen. Es ist das amerikanische Symbol, das aus mehreren aneinanderliegenden Schleifen besteht. Wir nutzen im Folgenden das der europäischen Norm entsprechende Schaltzeichen mit dem schwarzen Rechteck.}
	\silence{2}
\end{frame}

\begin{frame} \ftx{Bestimmung der Induktivität}
	\b{\begin{columns}
		\column[c]{0.5\textwidth}
		\f{width=0.5\textwidth}{height=0.7\textheight}{Spule_Ringkern1b}{{\bf Foto einer Ringkernspule.} Die ringförmige Bauweise bringt die Besonderheit mit sich, dass sich der magnetische Fluss vor allem innerhalb des Kerns bewegt, wodurch magnetische Störfelder reduziert werden.}

	\column[c]{0.5\textwidth}
	{\vspace{0pt}}%
	Bei einfachen geometrischen Strukturen kann die Induktivität über den magnetischen Widerstand bestimmt werden:
	{\vspace{1,5cm}}
	\begin{eq}
		L=\frac{N^2}{R_{\mathrm{m}}} \label{InduktivitaetFormel3}
	\end{eq}
	\end{columns}}
	\speech{InduktivitaetNaerungsformel}{1}{Wenn wir einfache geometrische Anordnungen betrachten, kann man die Induktivität mit einer Näherungsformel berechnen. Diese Vereinfachung funktioniert allerdings nur unter bestimmten Bedingungen. Das Magnetfeld sollte möglichst gleichmäßig sein, also über den gesamten Querschnitt hinweg dieselbe Stärke haben. Der Magnetkreis muss eine einfache und klar definierte Form besitzen, damit Länge, Fläche und Material bekannt sind. Zylinderspulen und Toroidspulen mit homogenem Material erfüllen oft diese Voraussetzungen. Zusätzlich, darf kein großer Luftspalt vorhanden sein, weil dieser das Magnetfeld stark beeinflusst. Das liegt daran, das Luft eine deutlich geringere magnetische Leitfähigkeit hat und somit der magnetische Widerstand merkbar steigt.}
	\silence{2}
	\speech{InduktivitaetNaerungsformel}{2}{Die Nährungsformel der Induktivität erfolgt über den Therm: N zum Quadrat geteilt durch R-M. Diese einfache Berechnung liefert einen guten ersten Anhaltspunkt, ist in der Praxis aber nie exakt. Der Grund liegt daran, dass Spulenwicklungen nie vollkommen gleichmäßig aufgebaut sind. Selbst bei maschinell gefertigten Spulen entstehen kleine Abweichungen, die das Magnetfeld beeinflussen. Wichtig ist, dass ihr euch merkt, dass diese Formel nur einen groben Schätzwert liefert und ihr sie nur anwenden solltet, wenn der Magnetkreis ein homogenes Feld besitzt und keinen oder nur einen geringen Luftspalt enthält.}
\end{frame}

\begin{frame} \ftx{Bestimmung der Induktivität}
	\s{Um die Induktivität einer Struktur zu bestimmen, gibt es zwei Möglichkeiten: Den Weg der Ermittlung über die Feldgrößen und den Weg über den magnetischen Widerstand.

	\begin{minipage}{0.49\textwidth}
		\vspace*{20pt}
		\textbf{Weg 1 über die Feldgrößen:}
		\begin{enumerate}%[label=\arabic*.]
			\item Annahme einer \\ Stromstärke $I$ $\rightarrow$ Durchflutung \par $\Theta = N\cdot I$
			\item Feldstärke berechnen \par $H=\frac{\Theta}{\ell_{\mathrm{m}}}$
			\item Flussdichte berechnen \par $B=\mu_{\mathrm{r}}\cdot\mu_0\cdot H$
			\item Magnetischen Fluss bestimmen \par $\varPhi = \iint_A \vec{B}\cdot \d \vec{A}$
			\item Induktivität \par $L=\frac{N\cdot \varPhi}{I}$\\
		\end{enumerate}
	\end{minipage}
	\begin{minipage}{0.49 \textwidth}
		\vspace*{-1.03cm}
		\textbf{Weg 2 über den magnetischen \\ Widerstand:}
		\begin{enumerate}%[label=\arabic*.]
			\item Magnetkreis in Teilabschnitte \\ zerlegen und magnetischen \\ Widerstand $R_{\mathrm{m}}$ berechnen \par $R_{\mathrm{m}}=\frac{\ell_{\mathrm{m}}}{\mu_{\mathrm{r}}\cdot\mu_0\cdot A}$
			\item Gesamtwiderstand berechnen\\(bei Reihenstruktur)\par $R_{m,\mathrm{ges}} = \sum R_{\mathrm{m}}$
			\item Induktivität \par $L=\frac{N^2}{R_{m,\mathrm{ges}}}$
		\end{enumerate}
	\end{minipage}
}
\b{
		\begin{minipage}{0.49\textwidth}
			\vspace*{20pt}
			\textbf{Weg 1 über die Feldgrößen:}
			\begin{enumerate}%[label=\arabic*.]
				\onslide<2->{\item Annahme einer \\ Stromstärke $I$ $\rightarrow$ Durchflutung \par $\Theta = N\cdot I$}
				\onslide<3->{\item Feldstärke berechnen \par $H=\frac{\Theta}{\ell_{\mathrm{m}}}$}
				\onslide<4->{\item Flussdichte berechnen \par $B=\mu_{\mathrm{r}}\cdot\mu_0\cdot H$}
				\onslide<5->{\item Magnetischen Fluss bestimmen \par $\varPhi = \iint_A \vec{B}\cdot \d \vec{A}$}
				\onslide<6->{\item Induktivität \par $L=\frac{N\cdot \varPhi}{I}$\\}
			\end{enumerate}
		\end{minipage}
		\begin{minipage}{0.49 \textwidth}
			\vspace*{-30pt}
			\onslide<7->{\textbf{Weg 2 über den magnetischen \\ Widerstand:}}
			\begin{enumerate}%[label=\arabic*.]
				\onslide<8->{\item Magnetkreis in Teilabschnitte \\ zerlegen und magnetischen \\ Widerstand $R_{\mathrm{m}}$ berechnen \par $R_{\mathrm{m}}=\frac{\ell_{\mathrm{m}}}{\mu_{\mathrm{r}}\cdot\mu_0\cdot A}$}
				\onslide<9->{\item Gesamtwiderstand berechnen\\(bei Reihenstruktur) \par $R_{m,\mathrm{ges}} = \sum R_{\mathrm{m}}$}
				\onslide<10->{\item Induktivität \par $L=\frac{N^2}{R_{m,\mathrm{ges}}}$}
			\end{enumerate}
		\end{minipage} 
}
	\speech{InduktivitaetRechenweg}{1}{Um die Induktivität einer unbekannten Struktur zu bestimmen, gibt es grundsätzlich zwei Wege. Ein Weg führt über die magnetischen Feldgrößen, wie die Durchflutung, Feldstärke, Flussdichte und magnetischen Fluss, die wir im Laufe dieses Kapitels kennengelernt haben. Der zweite Weg führt über den magnetischen Widerstand. Wir schauen uns jetzt beide Möglichkeiten näher an. Beginnen wir doch mit den Bestimmung über die magnetischen Feldgrößen.}
	\silence{2}
	\speech{InduktivitaetRechenweg}{2}{Lasst uns hierfür Schritt für Schritt vorgehen. Nehmen wir an, dass die Stromstärke in einer Spule bekannt ist. Wie wir zuvor erfahren haben, erzeugt dieser Strom eine magnetische Spannung, die wir als Durchflutung kennengelernt haben. Die Durchflutung Theta ergibt sich aus der Anzahl der Windungen multipliziert mit der Stromstärke.} 
	\silence{2}
	\speech{InduktivitaetRechenweg}{3}{Im zweiten Schritt können wir aus der Durchflutung die magnetische Feldstärke H berechnen. Hierzu teilen wir die Durchflutung Theta durch die mittlere Länge des magnetischen Kreises L M. Wie ihr euch sicher erinnert ist die mittlere Feldlinienlänge der angenäherte Weg, den die magnetischen Feldlinien im Inneren der Spule zurücklegen. Ist die magnetische Feldstärke H ermittelt können wir die magnetische Flussdichte berechnen.} 
	\silence{2}
	\speech{InduktivitaetRechenweg}{4}{Hierfür benötigen wir aber noch die Materialkonstante Mü-R. Ist diese bekannt, können wir die Flussdichte B über das Produkt von der materialspezifischen magnetischen Leitfähigkeit Mü-R, der magnetischen Feldkonstante mü-Null sowie der magnetischen Feldstärke H berechnen.} 
	\silence{2}
	\speech{InduktivitaetRechenweg}{5}{Aus der ermittelten magnetischen Flussdichte kann man dann wiederum den magnetischen Fluss Phi bestimmen, indem wir uns die Fläche anschauen, über die die Flussdichte wirkt. Hier müssen wir nun die magnetische Flussdichte B über die Fläche A integrieren. Das Flächenintegral wirkt sicher abschreckend. Verläuft jedoch die magnetische Flussdichte konstant über die Fläche A, lässt sich die magnetische Flussdichte vor das Integral ziehen und wir können den magnetischen Fluss Phi aus dem Produkt der Flussdichte und der Fläche A bestimmen.} 
	\silence{2}
	\speech{InduktivitaetRechenweg}{6}{Im letzten Schritt berechnen wir dann die Induktivität L über den therm der Spulenwindungen N multipliziert mit den magnetischen Fluss phi durch den Strom Ih. Wie ihr seht geht in dieser Ermittlung der Induktivität der Strom gleich zweimal ein. Einmal bei der Durchflutung und einmal bei der Berechnung der Induktivität. Somit taucht er bei der Durchflutung als Ursache des Magnetfelds auf und als Bezugsgröße der Induktivität selbst.}
	\silence{2}
	\speech{InduktivitaetRechenweg}{7}{Der zweite Weg eine Induktivität zu bestimmen, führt über den magnetischen Widerstand. Diesen Weg nutzen wir immer dann, wenn uns bestimmte Feldgrößen nicht bekannt sind – zum Beispiel, wenn wir die genaue Stromstärke nicht kennen. Welcher Weg besser passt, hängt also von der jeweiligen Aufgabe ab. Allerdings ist die Berechnung über den magnetischen Widerstand ungenauer, weil hier eine Näherungsformel verwendet wird. Des Weiteren können wir diese Formel nicht anwenden, wenn in der Anordung ein großer Luftspalt vorhanden ist. Kommen wir nun zu der Herangehensweise.} 
	\silence{2}
	\speech{InduktivitaetRechenweg}{8}{Wir zerlegen zuerst den Magnetkreis in einzelne Teilabschnitte, die jeweils einen konstanten magnetischen Widerstand haben. In einem Magnetkreis gibt es nämlich häufig verschiedene Bereiche – manche mit anderem Querschnitt, also einer anderen Fläche A, oder mit einem anderen Material, also einer unterschiedlichen relativen Permeabilität mü-R. Für jeden dieser Abschnitte berechnen wir den magnetischen Widerstand mit: R-M ist gleich L-M geteilt durch mü null mal mü R mal A. Dabei steht das L-M für die Länge des betrachteten Abschnitts, mü null für die Feldkonstante des Vakuums, mü-R für die relative Permeabilität des Materials, und A für den Querschnitt, durch den der magnetische Fluss im betrachteten Abschnitt läuft.} 
	\silence{2}
	\speech{InduktivitaetRechenweg}{9}{Haben wir die Widerstände der einzelnen Teilabschnitte berechnet, addieren wir sie anschließend zu einem Gesamtwiderstand R-M.} 
	\speech{InduktivitaetRechenweg}{10}{Mit diesem Gesamtwiderstand können wir schließlich die Induktivität mit der Nährungsformel bestimmen: L ist gleich N zum Quadrat durch R-M gesamt. So lässt sich also auch ohne genaue Kenntnis der Feldgrößen die Induktivität abschätzen. Im nächsten Video schauen wir uns dazu ein einfaches Rechenbeispiel an.}
	\silence{2}
\end{frame}

\newvideofile{Induktivitaet2}{Induktivität Teil 2}

\s{
	\begin{bsp}{Induktivität}{}
		Bei einer Ringkernspule mit Luftspalt soll der Einfluss des Luftspaltes $d$ auf die Induktivität $L$ untersucht werden. Die Spule hat $N=100$ Windungen. Die relative Permeabilität beträgt $\mu_{\mathrm{r}}=2000$. Die mittlere Eisenkernlänge beträgt $\ell_{\mathrm{m}}=5\,\mathrm{cm}$. Die Querschnittsfläche beträgt $A=1\,\mathrm{cm}^2$.

		\fo{width=0.4\textwidth, alt={XX}}{}{Ringspule_MagFeld_Luftspalt}{
			\put(5, 46){\color{current}$I$}
			\put(9, 0){\color{current}$I$}
			\put(60, 52){\color{flaeche}$A$}
			\put(68, 15){\color{magnetfeld}$\ell_{\mathrm{m}}$}
			\put(95, 55){\color{gray}$d$}
		}{{\bf Ringkernspule mit Luftspalt.} \label{AbbRingkern}}
		\begin{itemize}
			\item Der Luftspalt ist vorerst nicht vorhanden: $d=0$. Wie groß ist die Induktivität $L$?

			      \begin{align*}
				      R_{\mathrm{m}} & = \frac{\ell_{\mathrm{m}}}{\mu_{\mathrm{r}}\cdot\mu_0\cdot A}                                                                        \\
				      L          & = \frac{N^2}{R_{m}}                                                                                                          \\
				                 & = \frac{100^2\cdot 2000\cdot 1,256\cdot10^{-6}\,\frac{\mathrm{Vs}}{\mathrm{Am}}\cdot 1\cdot 10^{-4}\,\mathrm{m}^2}{0,05\,\mathrm{m}} \\
				      L          & = 0,05\,\frac{\mathrm{Vs}}{\mathrm{A}} = 50\,\mathrm{mH}
			      \end{align*}
			\item Der Luftspalt beträgt nun $d=1\,\mathrm{mm}$. Wie groß ist die Induktivität?
			      \begin{align*}
				      R_{\mathrm{m}} & = R_{m,\mathrm{Fe}} + R_{m,\mathrm{L}}                                                                                                 \\
				      L          & = \frac{\ell_{\mathrm{m}} - d}{\mu_{\mathrm{r}}\cdot\mu_0\cdot A} + \frac{d}{\mu_0\cdot A}                                                 \\
				      L          & =	\frac{0,05\,\mathrm{m} - 0,001\,\mathrm{m}}{2000\cdot 1,256\cdot10^{-6}\,\frac{\mathrm{Vs}}{\mathrm{Am}}\cdot 1\cdot 10^{-4}\,\mathrm{m}^2}
				      + \frac{0,001\,\mathrm{m}}{ 1,256\cdot10^{-6}\,\frac{\mathrm{Vs}}{\mathrm{Am}} \cdot 1\cdot 10^{-4}\,\mathrm{m}^2}                                      \\
				      L          & = 1,226\,\mathrm{mH}
			      \end{align*}
		\end{itemize}
	\end{bsp}

}
\b{
	\begin{frame}\ftx{Beispiel: Induktivität}
		Bei einer Spule mit Eisenkern und Luftspalt soll der Einfluss des Luftspaltes $d$ auf die Induktivität $L$ untersucht werden. Die Spule hat $N=100$ Windungen. Die relative Permeabilität beträgt $\mu_{\mathrm{r}}=2000$. Die mittlere Eisenkernlänge beträgt $\ell_{\mathrm{m}}=5\,\mathrm{cm}$. Die Querschnittsfläche beträgt $A=1\,\mathrm{cm}^2$.

		\foo{}{width=0.4\textwidth, alt={Das Bild zeigt einen ringförmigen Eisenkern, der von Drähten umwickelt ist. Rote Pfeile an den Drahtenden markieren die Richtung des elektrischen Stroms (I), der durch die Spule fließt. Neben dieser sogenannten Ringkernspule mit Luftspalt befindet sich in Grau die Darstellung des Luftspalts (d) und eine Fläche (A), die die Querschnittsfläche des Eisenkerns zeigt. Die Spule hat eine mittlere Feldlinienlänge (lm), die sich innerhalb des ringförmigen Eisenkerns befindet.}}{Ringspule_MagFeld_Luftspalt}{
			\put(5, 46){\color{current}$I$}
			\put(9, 0){\color{current}$I$}
			\put(60, 51){\color{flaeche}$A$}
			\put(68, 15){\color{magnetfeld}$\ell_{\mathrm{m}}$}
			\put(95, 55){\color{gray}$d$}
		}{}
	\end{frame}
	\begin{frame}\ftx{Beispiel: Induktivität}
		\begin{itemize}
			\item Der Luftspalt ist vorerst nicht vorhanden: $d=0$. Wie groß ist die Induktivität $L$?
		\pause
		\begin{align*}
			\onslide<1->{R_{\mathrm{m}} & = \frac{\ell_{\mathrm{m}}}{\mu_{\mathrm{r}}\cdot\mu_0\cdot A}} \\
			\onslide<2->{L & = \frac{N^2}{R_{\mathrm{m}}}} \\
			\onslide<3->{L & = \frac{100^2\cdot 2000\cdot 1,256\cdot10^{-6}\,\frac{\mathrm{Vs}}{\mathrm{Am}}\cdot 1\cdot 10^{-4}\,\mathrm{m}^2}{0,05\,\mathrm{m}}} \\
			\onslide<4->{L & = 0,05\,\tfrac{\mathrm{Vs}}{\mathrm{A}} = 50\,\mathrm{mH}}
		  \end{align*}


		\end{itemize}
	\end{frame}
	\begin{frame}\ftx{Beispiel: Induktivität}
		\begin{itemize}
			\item Der Luftspalt beträgt nun $d=1\,\mathrm{mm}$. Wie groß ist die Induktivität?
			\begin{align*}
				\onslide<2->{R_{\mathrm{m}} & = R_{m,\mathrm{Fe}} + R_{m,\mathrm{L}} \\
				& = \frac{\ell_{\mathrm{m}} - d}{\mu_{\mathrm{r}}\cdot\mu_0\cdot A} + \frac{d}{\mu_0\cdot A}} \\
				\onslide<3->{& = \frac{0,05\,\mathrm{m} - 0,001\,\mathrm{m}}{2000\cdot 1,256\cdot10^{-6}\,\frac{\mathrm{Vs}}{\mathrm{Am}}\cdot 1\cdot 10^{-4}\,\mathrm{m}^2} + \frac{0,001\,\mathrm{m}}{1,256\cdot10^{-6}\,\frac{\mathrm{Vs}}{\mathrm{Am}} \cdot 1\cdot 10^{-4}\,\mathrm{m}^2}} \\
				\onslide<4->{L & = 1,226\,\mathrm{mH}}
			\end{align*}
		\end{itemize}
	\end{frame}
}
\speech{InduktivitaetBeispiel}{1}{Nachdem wir die Grundlagen zur Berechnung einer Induktivität im vorherigen Video kennengelernt haben, wollen wir nun einen Blick auf ein einfaches Rechenbeispiel werfen. Um die Auswirkungen eines Luftspalts innerhalb einer Rinkernspule aufzuzeigen, betrachten wir hier eine Spule mit einem Eisenkern: einmal ohne und anschließend mit einem Luftspalt. In unseren Beispiel ist al Induktivität eine Spule mit Eisenkern mit einhundert Windungen gegeben. Die relative Permeabilität Mü R beträgt zweitausend. Die mittlere Länge des Eisenkerns beträgt L M gleich fünf Zentimeter und die Querschnittsfläche A ist gleich einem Quadratzentimeter.}
\speech{InduktivitaetBeispiel}{2}{Untersucht werden soll jetzt der Einfluss des Luftspalts D auf die Induktivität L. Hierfür betrachten wir zunächst den Aufbau ohne Luftspalt im Eisenkern.}
\speech{InduktivitaetBeispiel}{3}{Hier bietet es sich an, die Induktivität über den magnetischen Widerstand zu berechnen, weil dafür alle Werte vorliegen. Der magnetische Widerstand R M ergibt sich aus: R M ist gleich L M geteilt durch mü null mal Mü R mal A. Das A steht für die Querschnittsfläche. Diese Fläche ist bekannt und auch die mittlere Länge ist gegeben. Mit dem Widerstand allein sind wir aber noch nicht am Ziel angelangt, denn gesucht ist die Induktivität L. Dafür nutzen wir die Gleichung: L ist gleich N-Quadrat geteilt durch R M. Die Windungszahl N kennen wir ebenfalls, das heißt: Gesucht ist L, und rechts in der Gleichung stehen nur bekannte Größen – N, L M und A sowie die Permeabilitäten mü R und Mü-Null. Wir setzen also R M in die Formel für L ein und tragen anschließend Schritt für Schritt alle Zahlen ein.}
\speech{InduktivitaetBeispiel}{4}{Setzen wir nun alle Werte in die Formel ein, ergibt sich eine Induktivität L aus hundert Windungen zum Quadrat, mal zweitausend für die materialspezifische Permeabilität, mal der Permeabilität des Vakuums, also vier Pi mal zehn hoch minus sieben Voltsekunden pro Amperemeter, multipliziert mit der Querschnittsfläche von eins mal zehn hoch minus vier Quadratmetern, geteilt durch die mittlere Feldlinienlänge des Eisenkerns von null Komma null fünf Metern.}
\speech{InduktivitaetBeispiel}{5}{Rechnen wir alles aus, erhalten wir als Ergebnis null Komma null fünf Voltsekunden pro Ampere, die 50 milli Henry entsprechen.}
\speech{InduktivitaetBeispiel}{6}{Schauen wir uns jetzt bei der zweiten Rechnung an, was passiert, wenn wir in unseren Eisenkern einen kleinen Luftspalt von nur einen Millimeter einfügen. Wie verändert sich dadurch die Induktivität.}
\speech{InduktivitaetBeispiel}{7}{Also setzen wir wieder die Formel für den magnetischen Widerstand an: R M ist gleich L M geteilt durch Mü-Null mal Mü R mal A. Jetzt müssen wir aber den Luftspalt berücksichtigen. Dadurch zerfällt der Magnetkreis in zwei Teilbereiche, einmal in den magnetischen Widerstand des Eisens, den wir R M F-E nennen, und einmal in den magnetischen Widerstand der Luft, der hier R M-L genannt wird. In der Berechnung für den Widerstand im Eisenkern hat sich die Länge etwas verändert. Die wirksame Länge L M ist jetzt um den Luftspalt D verkürzt. Das heißt, wir verwenden in der Formel L M minus D. Das teilen wir wie zuvor durch mü null mal mü R mal A. Zusätzlich kommt jetzt der magnetische Widerstand der Luft hinzu. Das heißt, wir erweitern die Gleichung um den Anteil des Luftspalts. Den magnetischen Widerstand der Luft berechnen wir, indem wir die Länge des Luftspalts durch die Multiplikation der Permeabilität des Vakuums mit der Querschnittsfläche A teilen. Setzen wir nun die Werte ein. Für den magnetischen Widerstand des Eisens berechnen wir zuerst die Differenz aus der mittleren Länge von fünf Zentimetern minus dem Luftspalt von einem Millimeter. Diese Länge teilen wir durch die materialspezifische Permeabilität mü R gleich zweitausend, multipliziert mit der magnetischen Feldkonstanten mü null und mit der Querschnittsfläche von einem Quadratzentimeter. Anschließend errechnen wir den magnetischen Widerstand der Luft. Der ergibt sich aus dem Luftspalt von einem Millimeter, geteilt durch mü null und ebenfalls multipliziert mit der gleichen Querschnittsfläche von einem Quadratzentimeter. Zum Schluss addieren wir die beiden Werte.}
\speech{InduktivitaetBeispiel}{8}{Wenn wir das alles berechnen, erhalten wir eine Induktivität von ungefähr eins Komma zwei zwei sechs Milli Henry. Zum Vergleich: Ohne Luftspalt hatten wir eine Induktivität von etwa fünfzig Milli-Henry. Mit nur einem Millimeter Luftspalt, also etwa einem Fünfzigstel der Gesamtlänge, fällt die Induktivität auf nur noch eins Komma zwei Milli-Henry. Das zeigt deutlich, wie stark ein kleiner Luftspalt die magnetische Kopplung und damit die Induktivität reduziert. So, damit habt ihr nun ein solides Verständnis der Grundlagen der Induktion und der Induktivität. Ihr habt gesehen, wie sich durch die Änderung des Magnetfelds eine Spannung induzieren lässt, und wie dieses Wissen in der Praxis, in Form einer Induktivität, genutzt wird. Im nächsten Video schauen wir uns an, wie die Energie, die in einer Induktivität gespeichert ist, berechnet werden kann.}

\nomenclature[EH]{H}{Henry - Induktivität \nomunit{$\mathrm{H}=\frac{\mathrm{V}\cdot\mathrm{s}}{\mathrm{A}}=\frac{\mathrm{kg}\cdot\mathrm{m}^2}{\mathrm{s}^2\cdot\mathrm{A}^2}$}}
	\nomenclature[FL]{$L$}{Induktivität \nomunit{H}}